% THE MAP, third of three uses (\rulemap in short/style.tex). Ink = built; gold = proposed or modelled.
% No build: this slide stays up during questions.
\begin{frame}[t]{\ev{\tagB\,\tagP}\surf{THREE RULES}Contributions: three rules for doing work once.}
\vspace{.1cm}
\rulemap{\textbf{Cost rule}\\[4pt]{\color{Gold}\mbox{$(N{-}1)P > H+N(R{+}D)$}}\\[4pt]{\color{Gold}vs.\ a warm cache}}
{\textbf{Evidence contract}\\[4pt]Duplicate IDs rejected\\[4pt]{\color{Gold}Receipts, error weights}}
{\textbf{Acceptance gate}\\[4pt]Exact bytes, controller key\\[4pt]{\color{Gold}Promise tests, epoch fence}}{}
\claim{\centering\forktree\enspace Compute once. Count once. Publish once.}
\sources{Ink: built (publish check, duplicate-ID reducer). Gold: proposed or modelled. Factory runs: serial, nothing forked.}
\qadetail{The map from the rules slide a third time: each rule now has its contribution. The contribution claim is the integrated execution, evidence and acceptance contract; forking, caching, testing, review and provenance are established techniques, and no claim of invention, comprehensive priority or measured parallel speedup is made. Status by part: the cost rule is an additive resource model, priced against the warm alternative (restoring saved state versus rebuilding with a warm cache, both forms of reuse); H1 compares restoring a reached state with a warm cached template at N = 3, 6, 12 with batch makespan as the primary endpoint, which is not the same as total resource cost, and it has not started. The built reducer rejects duplicate declared observation IDs and carries lineage, but it does not abstain on dependence, and distinct IDs do not prove independence; typed receipts written outside the child and covariance weighting of shared errors are proposed. The factory runs exercised controller-held publication with no publish key in the worker, serially on Docker sandboxes, with nothing forked; the sha256-bound publish check is built; a required test for each promise, an external evaluator and the epoch fence are proposed (a TLA+ design exists, with no TLC run). A hash proves identity, not correctness. Two research questions remain: can an ancestry, fixtures and feedback manifest support a calibrated dependence model or a justified refusal of a precision claim, and can restore renew authority and reject stale publication across failure and uncertain external effects?}
\note{[1:00] My contributions, on the same map we started with. Compute once: a cost rule. Restore saved state only when the setup it avoids costs more than its overhead. Price it against a warm cache. Count once: an evidence contract. The reducer I built rejects duplicate observation IDs. Receipts kept outside the child and weights for shared errors are proposed. Publish once: an acceptance gate. The controller holds the publish key. Its hash-bound publish check is built. A test for every promise and the epoch fence are proposed. One hash threads all three. It keys the cache, binds the receipt, and names the published bytes. It proves identity, not correctness. Compute once. Count once. Publish once. Your questions, please.}
\end{frame}
